\documentclass[a4paper]{article}
\usepackage[margin=.8in]{geometry}
\usepackage{amsmath}
\usepackage{amssymb}
\usepackage{enumitem}
\usepackage[cache=true]{minted}
\usepackage{fancyhdr}
\newcommand{\julia}[1]{\mintinline{julia}{#1}}
\setminted{fontsize=\small,autogobble}

\pagestyle{fancy}
\fancyhf{}
\lhead{Name: }
\chead{Graded quiz 2 -- Introduction to Julia}
\rhead{\today}
\cfoot{\thepage}

% Define \showsolutions when compiling to append the instructor answer key.
\begin{document}

True or false? (unless otherwise specified.)
Assume any explicitly named package is installed.
An expression that raises an error does not evaluate to \julia{true}.

\begin{enumerate}[itemsep=8pt]
    \item
        At an empty Julia REPL prompt, pressing \julia{?} enters package mode,
        where packages can be installed.

    \item
        After \julia{import Plots}, the command
        \julia{Plots.plot(x -> sin(x) + cos(x), 0, 2pi)} plots the function
        $x \mapsto \sin(x) + \cos(x)$ on $[0,2\pi]$.

    \item
        In a Jupyter notebook using Julia, a variable assigned in one cell
        cannot be used in another cell, even if both cells are executed
        in order using the same kernel.

    \item
        The expression \julia{typeof(1 + 2.0) == Float64} evaluates to \julia{true}.

    \item
        The command \julia{collect(2:2:10)} returns the vector
        \julia{[2, 4, 6, 8, 10, 12]}.

    \item
        If \julia{A = [1 2 3; 4 5 6]}, then \julia{A[1, :]} is a
        $1 \times 3$ matrix, rather than a vector.

    \item
        The following code assigns the vector \julia{[2, 5]} to \julia{w}.
        \begin{minted}{julia}
            v = [-3, 2, 0, 5, -1]
            w = v[v .> 0]
        \end{minted}

    \item
        The following code assigns the vector \julia{[5.0, 13.0]} to \julia{a}.
        \begin{minted}{julia}
            hypotenuse(x, y) = sqrt(x^2 + y^2)
            a = hypotenuse.([3, 5], [4, 12])
        \end{minted}

    \item
        The expression \julia{sum(n -> n^2 - 1, [1, 2, 3])} returns \julia{14}.

    \item
        The comprehension \julia{[n^2 for n in 1:5 if isodd(n)]}
        returns the vector \julia{[1, 9, 25]}.
\end{enumerate}

\newpage
\begin{enumerate}[resume,itemsep=8pt]
    \item
        After the following definitions, \julia{f(f(1))} returns \julia{2.0}.
        \begin{minted}{julia}
            f(x::Float64) = 1
            f(x::Int) = 2.0
        \end{minted}

    \item
        The call \julia{odd_sum(5)} returns \julia{9}.
        \begin{minted}{julia}
            function odd_sum(n)
                s = 0
                for k in 1:n
                    if isodd(k)
                        s += k
                    end
                end
                return s
            end
        \end{minted}

    \item
        After the definition \julia{p(n) = n == 0 ? 1 : n * p(n - 1)},
        the call \julia{p(4)} returns \julia{24}.

    \item
        After the definition \julia{shift(x; a=1) = x + a}, both
        \julia{shift(3, 2)} and \julia{shift(3; a=2)} return \julia{5}.

    \item
        After executing the following code, \julia{v} is still \julia{[1, 2, 3]}.
        \begin{minted}{julia}
            v = [1, 2, 3]
            w = v
            w[1] = 10
        \end{minted}
\end{enumerate}

\begin{center}
    \textbf{Answer grid}\par\smallskip
    \small
    \renewcommand{\arraystretch}{1.25}
    \begin{tabular}{|c|*{15}{c|}}
        \hline
        & 1 & 2 & 3 & 4 & 5 & 6 & 7 & 8 & 9 & 10 & 11 & 12 & 13 & 14 & 15 \\
        \hline
        True & $\square$ & $\square$ & $\square$ & $\square$ & $\square$ & $\square$ & $\square$ & $\square$ & $\square$ & $\square$ & $\square$ & $\square$ & $\square$ & $\square$ & $\square$ \\
        False & $\square$ & $\square$ & $\square$ & $\square$ & $\square$ & $\square$ & $\square$ & $\square$ & $\square$ & $\square$ & $\square$ & $\square$ & $\square$ & $\square$ & $\square$ \\
        \hline
    \end{tabular}
\end{center}

\noindent\textbf{Bonus (2 points).}
Give the values of \julia{a} and \julia{b} after the following code,
and briefly explain why the results differ.
\begin{minted}{julia}
    v = [[1, 2, 3], [2, 3, 4], [4, 5, 6]]
    a = sum.(v)
    b = sum(v)
\end{minted}
\noindent\textbf{Answer and explanation:}
\vspace{2cm}

\ifdefined\showsolutions
\newpage
\lhead{Instructor copy}
\section*{Answer key}
\noindent Questions 1--15: 1 point for a correct answer, 0 otherwise.

\begin{enumerate}[itemsep=8pt]
    \item \textbf{False.} \julia{?} enters help mode; \julia{]} enters package mode.

    \item \textbf{True.} \julia{import Plots} makes the module name available.
        The qualified call \julia{Plots.plot} accepts the anonymous function
        and the two interval endpoints.

    \item \textbf{False.} Cells share the kernel's state. A variable is available
        after its defining cell has been executed, unless it is subsequently
        removed or the kernel is restarted.

    \item \textbf{True.} The integer is promoted for this addition;
        the result is \julia{3.0}, of type \julia{Float64}.

    \item \textbf{False.} The range stops at 10, so the result is
        \julia{[2, 4, 6, 8, 10]}.

    \item \textbf{False.} Scalar indexing of the first dimension drops that
        dimension: \julia{A[1, :]} is the vector \julia{[1, 2, 3]}.
        Use \julia{A[1:1, :]} to retain a $1 \times 3$ matrix.

    \item \textbf{True.} The Boolean mask selects exactly the strictly positive
        entries, preserving their order.

    \item \textbf{True.} Broadcasting pairs 3 with 4, and 5 with 12.
        Their hypotenuses are 5 and 13, respectively.

    \item \textbf{False.} The mapped values are 0, 3, and 8,
        whose sum is 11. The expression \julia{sum(n -> n^2, [1, 2, 3])}
        would return 14.

    \item \textbf{True.} The filter retains 1, 3, and 5 before squaring them.

    \item \textbf{False.} \julia{f(1)} returns \julia{2.0}; the outer call
        therefore uses the \julia{Float64} method and returns \julia{1}.

    \item \textbf{True.} The accumulator adds $1+3+5=9$.

    \item \textbf{True.} The recursion computes
        $4\times3\times2\times1\times1=24$ and stops at \julia{p(0)}.

    \item \textbf{False.} \julia{a} is a keyword argument.
        The call \julia{shift(3, 2)} raises a \julia{MethodError};
        \julia{shift(3; a=2)} returns 5.

    \item \textbf{False.} \julia{w = v} does not copy the array.
        Both names refer to the same array, so \julia{v} becomes
        \julia{[10, 2, 3]}. Using \julia{w = copy(v)} would avoid this mutation
        of \julia{v}.
\end{enumerate}

\noindent\textbf{Bonus.}
\julia{a = [6, 9, 15]} and \julia{b = [7, 10, 13]}.
Broadcasting \julia{sum} applies it separately to each inner vector,
whereas \julia{sum(v)} adds the three vectors componentwise.
Award 1 point for each correct vector with its explanation
(0.5 for the vector and 0.5 for the explanation).
\fi

\end{document}
